\documentclass[10pt,a4paper]{amsart}
\usepackage{amsmath,amssymb,amsthm,graphicx}
% Retained preamble for the Stacks Project chapter {CH} (shared project preamble).
% The chapter file pinned-original-source.tex only carries \input{{preamble}} on its first line,
% so no preamble text lives inside the chapter itself.  The Stacks Project uses one single shared
% preamble for every chapter: the snapshot stages it three times (sources/stacks-chow/preamble.tex,
% sources/stacks-etale-cohomology/preamble.tex and focused-seeds/AGP/sources/stacks-snapshot-20261011/
% preamble.tex) and those copies are identical in size.  The block below is copied mechanically,
% line for line, from sources/stacks-chow/preamble.tex lines 14-18 and 154-217: the Xy-pic setup and
% the native theorem declarations together with the mathematical macros.  Nothing was added, removed
% or redefined; the wrapper class and the amsmath/amssymb/amsthm/graphicx base packages are supplied
% by the assembler.
% For commutative diagrams we use Xy-pic
% Chapter: stacks-cohomology.tex -- Cohomology of Sheaves.
\usepackage[all]{xy}

% We use 2cell for 2-commutative diagrams.
\xyoption{2cell}

% Theorem environments.
%
\theoremstyle{plain}
\newtheorem{theorem}[subsection]{Theorem}
\newtheorem{proposition}[subsection]{Proposition}
\newtheorem{lemma}[subsection]{Lemma}

\theoremstyle{definition}
\newtheorem{definition}[subsection]{Definition}
\newtheorem{example}[subsection]{Example}
\newtheorem{exercise}[subsection]{Exercise}
\newtheorem{situation}[subsection]{Situation}

\theoremstyle{remark}
\newtheorem{remark}[subsection]{Remark}
\newtheorem{remarks}[subsection]{Remarks}

\numberwithin{equation}{subsection}

% Macros
%
\def\lim{\mathop{\mathrm{lim}}\nolimits}
\def\colim{\mathop{\mathrm{colim}}\nolimits}
\def\Spec{\mathop{\mathrm{Spec}}}
\def\Hom{\mathop{\mathrm{Hom}}\nolimits}
\def\Ext{\mathop{\mathrm{Ext}}\nolimits}
\def\SheafHom{\mathop{\mathcal{H}\!\mathit{om}}\nolimits}
\def\SheafExt{\mathop{\mathcal{E}\!\mathit{xt}}\nolimits}
\def\Sch{\mathit{Sch}}
\def\Mor{\mathop{\mathrm{Mor}}\nolimits}
\def\Ob{\mathop{\mathrm{Ob}}\nolimits}
\def\Sh{\mathop{\mathit{Sh}}\nolimits}
\def\NL{\mathop{N\!L}\nolimits}
\def\CH{\mathop{\mathrm{CH}}\nolimits}
\def\proetale{{pro\text{-}\acute{e}tale}}
\def\etale{{\acute{e}tale}}
\def\QCoh{\mathit{QCoh}}
\def\Ker{\mathop{\mathrm{Ker}}}
\def\Im{\mathop{\mathrm{Im}}}
\def\Coker{\mathop{\mathrm{Coker}}}
\def\Coim{\mathop{\mathrm{Coim}}}

% Boxtimes
%
\DeclareMathSymbol{\boxtimes}{\mathbin}{AMSa}{"02}

%
% Macros for moduli stacks/spaces
%
\def\QCohstack{\mathcal{QC}\!\mathit{oh}}
\def\Cohstack{\mathcal{C}\!\mathit{oh}}
\def\Spacesstack{\mathcal{S}\!\mathit{paces}}
\def\Quotfunctor{\mathrm{Quot}}
\def\Hilbfunctor{\mathrm{Hilb}}
\def\Curvesstack{\mathcal{C}\!\mathit{urves}}
\def\Polarizedstack{\mathcal{P}\!\mathit{olarized}}
\def\Complexesstack{\mathcal{C}\!\mathit{omplexes}}
% \Pic is the operator that assigns to X its picard group, usage \Pic(X)
% \Picardstack_{X/B} denotes the Picard stack of X over B
% \Picardfunctor_{X/B} denotes the Picard functor of X over B
\def\Pic{\mathop{\mathrm{Pic}}\nolimits}
\def\Picardstack{\mathcal{P}\!\mathit{ic}}
\def\Picardfunctor{\mathrm{Pic}}
\def\Deformationcategory{\mathcal{D}\!\mathit{ef}}
\usepackage{cleveref}
\crefname{originallemma}{original lemma}{original lemmas}
\Crefname{originallemma}{Original lemma}{Original lemmas}
\makeatletter
\newlabel{lemma-cohomology-with-support-sheaf-on-support}{{0A3B}{source}{}{sourcepointer-1}{}}
\newlabel{lemma-cohomology-with-support-sheaf-on-support@cref}{{[originallemma][0][]0A3B}{[1][0][]source}{}{}{}}
\makeatother
\crefname{originalreference}{original reference}{original references}
\Crefname{originalreference}{Original reference}{Original references}
\makeatletter
\newlabel{modules-lemma-exactness-pushforward-pullback}{{01AJ}{source}{}{sourcepointer-2}{}}
\newlabel{modules-lemma-exactness-pushforward-pullback@cref}{{[originalreference][0][]01AJ}{[1][0][]source}{}{}{}}
\makeatother
\crefname{originalreference}{original reference}{original references}
\Crefname{originalreference}{Original reference}{Original references}
\makeatletter
\newlabel{modules-lemma-i-star-exact}{{01AX}{source}{}{sourcepointer-3}{}}
\newlabel{modules-lemma-i-star-exact@cref}{{[originalreference][0][]01AX}{[1][0][]source}{}{}{}}
\makeatother
\title{The derived category of modules supported on a closed subset is equivalent to the derived category of the closed subspace, with the functor as right adjoint reconstruction}
\author{The Stacks Project Authors}
\date{stacks-project:cohomology.tex:local-snapshot-20261011; GNU FDL 1.2 or later}
\begin{document}
\maketitle
\noindent\textit{Editorial note.} All delivered mathematical text is the author's own, copied verbatim from the pinned Stacks Project chapter file cohomology.tex; no statement or proof step was rewritten, reordered or removed.\par\medskip
\noindent\textit{Editorial note.} The counted span cites three cross-chapter labels, each recorded in assembly.external\_references with its canonical tag; none replaces a load-bearing step, which are all inside the span.\par\medskip
\noindent\textit{Editorial note.} The proof cites no bibliography item inside the counted span.\par\medskip
\noindent\textit{Allowed standard prerequisites (editorial).}\begin{itemize}
\item Exactness of pushforward and pullback along a closed immersion. Stacks Project, Modules on Sites, Lemmas modules-lemma-exactness-pushforward-pullback (tag 01AJ) and modules-lemma-i-star-exact (tag 01AX). Supplies exactness of i\_* and i\textasciicircum{}* = i\textasciicircum{}\{-1\}, which drives full faithfulness of i\_*.
\item Cohomology with support on a sheaf supported in the closed subset. Stacks Project, Cohomology of Sheaves, Lemma lemma-cohomology-with-support-sheaf-on-support (tag 0A3B). Identifies R-H\_Z with i\textasciicircum{}\{-1\} on D\_Z(O\_X) and gives part (3).
\end{itemize}\medskip
\ifdefined\hypertarget\hypertarget{sourcepointer-1}{}\fi
\noindent\textit{Source reference.} Original lemma 0A3B: sources/stacks-chow/tags.txt, source lines 12734--12734. Cross-chapter standard prerequisite cited by the author inside the counted span; it is used as a pointer and its own proof is not part of this item.\par\medskip
\ifdefined\hypertarget\hypertarget{sourcepointer-2}{}\fi
\noindent\textit{Source reference.} Original reference 01AJ: sources/stacks-chow/tags.txt, source lines 1692--1692. Cross-chapter standard prerequisite cited by the author inside the counted span; it is used as a pointer and its own proof is not part of this item.\par\medskip
\ifdefined\hypertarget\hypertarget{sourcepointer-3}{}\fi
\noindent\textit{Source reference.} Original reference 01AX: sources/stacks-chow/tags.txt, source lines 1706--1706. Cross-chapter standard prerequisite cited by the author inside the counted span; it is used as a pointer and its own proof is not part of this item.\par\medskip
\begin{lemma}
\label{lemma-complexes-with-support-on-closed}
Let $(X, \mathcal{O}_X)$ be a ringed space. Let $i : Z \to X$ be the
inclusion of a closed subset.
\begin{enumerate}
\item For $K$ in $D(\mathcal{O}_X|_Z)$ we have $i_*K$ in $D_Z(\mathcal{O}_X)$.
\item The functor $i_* : D(\mathcal{O}_X|_Z) \to D_Z(\mathcal{O}_X)$
is an equivalence with quasi-inverse
$i^{-1}|_{D_Z(\mathcal{O}_X)} = R\mathcal{H}_Z|_{D_Z(\mathcal{O}_X)}$.
\item The functor
$i_* \circ R\mathcal{H}_Z : D(\mathcal{O}_X) \to D_Z(\mathcal{O}_X)$
is right adjoint to the inclusion functor
$D_Z(\mathcal{O}_X) \to D(\mathcal{O}_X)$.
\end{enumerate}
\end{lemma}

\begin{proof}
Part (1) is immediate from the definitions. Part (3) is a formal
consequence of part (2) and
Lemma \ref{lemma-cohomology-with-support-sheaf-on-support}.
In the rest of the proof we prove part (2).

\medskip\noindent
Let us think of $i$ as the morphism of ringed spaces
$i : (Z, \mathcal{O}_X|_Z) \to (X, \mathcal{O}_X)$.
Recall that $i^*$ and $i_*$ is an adjoint pair of functors.
Since $i$ is a closed immersion, $i_*$ is exact.
Since $i^{-1}\mathcal{O}_X = \mathcal{O}_X|_Z$ is the structure
sheaf of $(Z, \mathcal{O}_X|_Z)$ we see that $i^* = i^{-1}$
is exact and we see that that $i^*i_* = i^{-1}i_*$
is isomorphic to the identify functor. See
Modules, Lemmas \ref{modules-lemma-exactness-pushforward-pullback} and
\ref{modules-lemma-i-star-exact}. Thus
$i_* : D(\mathcal{O}_X|_Z) \to D_Z(\mathcal{O}_X)$
is fully faithful and $i^{-1}$ determines
a left inverse. On the other hand, suppose that $K$ is an object of
$D_Z(\mathcal{O}_X)$ and consider the adjunction map
$K \to i_*i^{-1}K$.
Using exactness of $i_*$ and $i^{-1}$ this induces the adjunction maps
$H^n(K) \to i_*i^{-1}H^n(K)$ on cohomology sheaves. Since these cohomology
sheaves are supported on $Z$ we see these adjunction maps are isomorphisms
and we conclude that $i_* : D(\mathcal{O}_X|_Z) \to D_Z(\mathcal{O}_X)$
is an equivalence.

\medskip\noindent
To finish the proof it suffices to show that $R\mathcal{H}_Z(K) = i^{-1}K$ if
$K$ is an object of $D_Z(\mathcal{O}_X)$. To do this we can use that
$K = i_*i^{-1}K$ as we've just proved this is the case. Then
Lemma \ref{lemma-cohomology-with-support-sheaf-on-support}
tells us what we want.
\end{proof}

\end{document}
